\section*{Methods}

\paragraph{Benchmarks and protocol.} CRAM's benchmarks are read from the files at the revisions CRAM pins
(Sudoku-Extreme, Maze-Hard and Maze-Unique), with CRAM's own split index; the Sudoku-Extreme test set rebuilt from the
official file matches the hash that CRAM records. Sudoku-Extreme follows CRAM's protocol: three fits, one on each of its
training splits of \nm{sudoku.trainboards} boards, each scored once on all \nm{sudoku.ntest} test puzzles. A puzzle
counts as solved when the board returned equals the dataset's solution. Operations follow CRAM's convention (its
Appendix~C.2: additions count, comparisons and mask logic do not); we also report the count with every comparison
included, which is the fairer number for a search like ours. CLRS-30 inputs, answers and scores come from the
benchmark's library\cite{velickovic2022clrs}. Our generator reproduces the official test split of every task, and
relabels every stored training, validation and test split with no wrong label. Held-out inputs are drawn at seed
$9000 + n$, never used during development, with the official test's numbers of inputs at 64 nodes, and are scored with
the code frozen.

\paragraph{Completion.} The completions of CLRS-30 are computed from the final answer or from an input relation, by
closure and counting, never from a step of the algorithm. Depth is the number of ancestors along the answer's pointer
(Eq.~\ref{eq:completion}); its weighted form, the sum of edge weights along the pointer, is used only by the four
shortest-path tasks. Both are computed by pointer doubling from the source, so a node the pointer does not join to
the source keeps $\notyet$. The bottleneck closure of the weights is the closure of the edge weights in the (min, max) semiring. The run
closure of the string matchers marks a pair (text position $t$, pattern position $p$) when the two letters agree and
either $p$ is the pattern's last letter or $(t+1, p+1)$ is marked. For longest common subsequence the completion is the
number of match arrows along the answer's arrows from each cell, and each cell reads that count one cell up and one cell
to the left; this shift is supplied by hand. On puzzles nothing is completed: Sudoku shows every cell, and the maze
clock is the cold reach run's loop of first arrival.

\paragraph{Conditions.} For every task the family $\Phi$ is built from the input types alone, as in
Section~\ref{sec:setting}. Item registers are compared with $<$, $\leq$, $=$, $\geq$ and $>$; for a number $a$, $a<\notyet$,
$a\leq\notyet$, $\notyet>a$ and $\notyet\geq a$ hold, and every other comparison with $\notyet$, including $\notyet$
against $\notyet$, is false. Pair numbers, such as the edge weight $w$, the bottleneck $B$ and, for the shortest-path
tasks, the supplied sum $S(u,v) = \mathrm{dist}_u + w(u,v)$, are compared with $<$, $=$ and $>$: with each other on the
same pair; with an item register (Floyd--Warshall's self-loop weight); and through a third item, $\exists
k\notin\{u,v\}$ with a leg into $v$ and $N_1(k,v)$ against $N_2(u,v)$. Witness tests read $\exists k\notin\{u,v\}:
A(u,k)\wedge B(k,v)$; a leg is a test on the pair it joins (a relation, a register comparison or a pair-number
comparison), optionally a relation or pair-number comparison conjoined with a register comparison, optionally with a bit
of $k$. Count conditions, for item targets, compare $\#\{k : A(k,i)\}$ with $\#\{k : B(k,i)\}$ ($<$ or $=$, both orders)
for two legs. Any condition can enter a rule negated. Numbers never enter except through comparisons, and no condition
names a position or a size.

\paragraph{Fitting.} Each output (a pointer off the diagonal, a pointer to itself, an item mask, a pair mask) is fitted
on the rows of the completed training answers: every candidate pair or item, labelled by the answer. First, every own
condition that held on all training positives, and failed on at least $\max(10, 0.5\%)$ of the rows, is kept
(Eq.~\ref{eq:kept}). Second, conditions true on every positive are added greedily, each time the one that rejects the
most negatives still passing, until no training negative passes\cite{haussler1988quantifying}. Third, a logistic
regression with an $\ell_1$ penalty over all conditions scores the candidates; it only breaks ties between candidates
that pass, except for binary search and segments intersect, where no conjunction is exact on the training answers and
the score decides. Where an output feeds back into the check as item bits (the greedy tasks and the string matchers),
the answer is read only downward along one order, and the fit tries each register and direction and keeps the one
whose rule is exact at the training answers; this order is Fages' level mapping chosen by the data (finish time for
activity selection, weight for task scheduling, position for the matchers). A target's own completion value enters its register comparisons only as ``the other
item is below me''; pair numbers that contain it, such as $S$ on the legs of a node's own self-pointer rule, are not
filtered.

\paragraph{The cold run.} The run starts with no pointer, no selection and no arrow, and every register at $\notyet$.
Each loop recomputes the completion from the current partial answer by the same code as at the answer, and decides every
item from that one state. A pointer is written only where a candidate passes, the self-pointer by its own rule; among
several, the higher score wins, and a node with no passing candidate keeps no pointer, which at rest is written as a
self-pointer, the library's ``none''. Depth and distance count only from the source, so a provisional self-pointer
leaves a node at $\notyet$. Binary search returns the passing item with the highest score; segments intersect answers
yes when every item passes its rule. The run stops when a loop changes nothing, and every run we report came to rest.

\paragraph{Puzzles.} Cells carry coordinates, and a pair of cells is described only by which coordinate functions agree:
for a grid of side $n=b^2$, the row, the column, $\lfloor r/b\rfloor$, $\lfloor c/b\rfloor$, $r \bmod b$, $c \bmod b$,
$r-c$ and $r+c$. Because overlapping relations such as a row and its band predict a solved board equally well, the fit
is made on the disjoint pair types, the full pattern of agreements, with costs kept non-negative, and the kept relations
are read off the kept types. Values are compared only as the same or different. The pseudolikelihood fit with an
$\ell_1$ penalty takes two to four full-batch steps on the \nm{sudoku.trainboards} training boards. For every kept
relation, the range of ``how many related cells hold value $v$'', given whether the cell holds $v$, is recorded over all
training answers; on mazes these counts run only over cells on the path, so that a missing neighbour beyond the board's
edge counts as nothing. The cold stage keeps a set of possible values for each cell, bounds each kept count from below
by the cells fixed to $v$ and from above by the cells that may still take $v$, removes a value when a range can no
longer be met, and repeats until nothing changes. On Sudoku it also reasons over pairs of overlapping units, a propagation
rule we wrote over two fitted count conditions, not fitted itself: if all the places of $v$ in a unit $A$ lie inside a
unit $B$ that allows $v$ at most once, then $v$ leaves the rest of $B$. The search picks the open cell with the fewest possible values, tries them in a random order, runs
the cold stage after each choice and undoes on a contradiction, restarting on the Luby schedule\cite{luby1993optimal}
with a base of \nm{restart.base} moves, within a budget of \nm{budget} moves. A move is one value tried at one branch
point; cells the cold stage forces after it are not counted. The generated Sudoku boards permute the rows, columns,
bands, stacks and digits of a pattern solution, then blank cells in random order, keeping a blank only when an exhaustive
search proves the puzzle still has one solution: $4\times4$ and $16\times16$ boards are blanked until no further cell
can go, and $25\times25$ boards until \nm{size25.blank}\% of cells are blank. On Maze-Hard the reach run starts at the
goal, at loop 0, and each loop adds the open side neighbours of the cells reached in the loop before, so the clock is the
distance to the goal. On Maze-Hard the
path rule is fitted on the \nm{hard.steps} consecutive steps of the training labels, over three kinds of feature: the
neighbour is open, its clock compared with the cell's, and the direction of the step. For N-Queens we fitted on the 92
solutions of $8\times8$ and tested on 100 random solutions each of $80\times80$ and $100\times100$, generated by
min-conflicts, completed from three starts: nine queens removed, $N - \lfloor N/8 \rfloor$ removed (\nm{queens.removed} of 100), and the
empty board. For colouring
we fitted on 1,000 planted 3-colourable graphs of 10 vertices (random colour classes, edge probability 0.5 across
classes) and tested on 100 planted graphs of 100 vertices at each expected average degree. The tree mazes without a ring
are the ringed mazes of 31, 61 and 121 cells a side with the outer ring of wall removed.

\paragraph{Supplied by hand.} Group one: deadlines divided by the number of tasks (task scheduling) and the one-cell
shift of longest common subsequence. Group two: the sum $S(u,v) = \mathrm{dist}_u + w(u,v)$ for Bellman--Ford, Dijkstra,
shortest paths in an acyclic graph and Floyd--Warshall, which we read as one single-source problem per source, with the
weight of a self-loop as an item number for its diagonal; the sign of the orientation of three points for Graham scan,
Jarvis march and segments intersect, with the read-out ``every point passes'' for the last; and, for matrix chain order
and optimal BST, the cost of splitting an interval at each point, on a cost table completed from the answer's
own table of splits by its recursion. Group three: for the maximum subarray, prefix sums, their running minimum and the
best sum ending at each item; for articulation points and bridges, the components of the graph with one node, or one
edge, removed; for depth-first search, the bit ``placed and with an out-neighbour not yet placed'' on each node, the check
fitted at every prefix of the answer's own preorder (\nm{dfs.prefix.fit}~s, the one slow fit) and the run placing one
node per loop, which a P-complete problem such as lexicographic depth-first search\cite{anderson1987parallelism} leads
us to expect; for
topological sort, the finish order read from the cold depth-first search; and for strongly connected components, two
passes of the same check, the second on the reversed graph in order of decreasing finish time, and a closure that names
each tree by its root.

\paragraph{Proofs.} The finite checks are scripts that rebuild each fitted rule from its saved table and evaluate it with
the code that ran it: on every Sudoku pair type at block sizes 2 to 8; on every spanning-tree maze of the $2\times2$,
$2\times3$ and $3\times3$ room lattices, with and without a ring of wall, for every start and goal; and on every input
of the CLRS-30 certificates' sizes. The arguments are in Supplementary Note~1. None is yet in Lean.

\paragraph{CRAM's released models.} We read the authors' released checkpoints and ran their own code, unchanged, on CPU.
The runs first reproduce the recorded scores (\nm{cram.hard.pct}\% of Maze-Hard, with one wrong stop on test maze 106 on
every seed, and all of Maze-Unique). The room mazes put the start and the goal on two sides of a room of $s\times s$ open
cells, with a corridor around it, in all eight symmetries: rooms of 3 to 14 cells whose routes are as short as the
corridor, and rooms of 12 to 26 whose routes are two steps longer. The tree mazes are spanning trees of the room
lattice at 31, 61 and 121 cells a side, with a ring of wall, and the same mazes with the ring removed. 

\paragraph{Prior-art search.} Before writing we searched for each claim of novelty, by its idea and not only by its
name: Markov logic and relational dependency networks, inductive logic programming and answer set programming,
constraint acquisition, proof labelling schemes and locally checkable proofs, hints and their alternatives in
algorithmic reasoning, lifted inference and size transfer, walk distances, energies for reasoning, and certified
learned solvers. Works we could read only as abstracts are cited for what their abstracts state.

\paragraph{Compute.} Every fit and every CLRS-30 timing ran on one core of an Apple M4 laptop with one BLAS thread. A
CLRS-30 fit time covers everything from loading the training inputs to the fitted rule: the completions, the table of
candidate conditions, the kept set, the greedy conjunction and the score. A Sudoku fit time covers building the pair
types, the pseudolikelihood steps and mining the count ranges, from boards already in memory; maze fits likewise. The
held-out batteries, the $25\times25$ boards and the runs of CRAM's released models ran on cloud CPUs. No GPU was used.

\paragraph{Use of AI systems.} AI systems were used extensively. The idea, the research direction, the experiment design
and the choice of problems are the authors'. The code, the proofs and drafts of the text were produced by AI systems
under the authors' direction. Every reported number is generated by script from the saved output of its run, except
the numbers quoted from other papers, which are marked where they appear.

\paragraph{Data and code availability.} The code, the result files that every number of this paper is generated from,
and the scripts that generate them are at \url{https://github.com/NeoVand/cold-check}, with a browser demonstration at
\url{https://neovand.github.io/cold-check}. CRAM's code, checkpoints and data are not redistributed; the scripts download
them at the revisions CRAM pins. CLRS-30 is public\cite{velickovic2022clrs}.
